% Part: first-order-logic
% Chapter: axiomatic-deduction
% Section: proof-theoretic-notions

\documentclass[../../../include/open-logic-section]{subfiles}

\begin{document}

\iftag{FOL}
      {\olfileid{fol}{axd}{ptn}}
      {\olfileid{pl}{axd}{ptn}}

\olsection{Bewijstheoretische begrippen}

\begin{explain}
Zoals we een aantal belangrijke semantische begrippen hebben
gedefinieerd (\iftag{FOL}{geldigheid}{tautologie}, gevolg en
vervulbaarheid), definiëren we nu de overeenkomstige
\emph{bewijstheoretische begrippen}. Deze worden niet gedefinieerd aan
de hand van het waar zijn van !!{sentence}s in !!{structure}s, maar aan
de hand van de !!{derivability} of !!{nonderivability} van bepaalde
formules. Het was een belangrijke ontdekking dat deze begrippen
samenvallen. Dat ze samenvallen, is de inhoud van de
\emph{correctheids-} en \emph{volledigheidsstellingen}.
\end{explain}

\begin{defn}[!!^{derivability}]
!!^a{formula}~$!A$ is \emph{!!{derivable} uit} $\Gamma$, genoteerd als
$\Gamma \Proves !A$, als er !!a{derivation} uit~$\Gamma$ bestaat die
eindigt met~$!A$.
\end{defn}

\begin{defn}[Stellingen]
!!^a{formula}~$!A$ is een \emph{stelling} als er !!a{derivation} van
$!A$ uit de lege verzameling bestaat. We schrijven $\Proves !A$ als
$!A$ een stelling is en $\Proves/ !A$ als dat niet zo is.
\end{defn}

\begin{defn}[Consistentie]
Een verzameling !!{formula}s~$\Gamma$ is \emph{consistent} dan en
slechts dan als $\Gamma\Proves/ \lfalse$; anders is zij
\emph{inconsistent}.
\end{defn}

\begin{prop}[Reflexiviteit]
\ollabel{prop:reflexivity}
Als $!A \in \Gamma$, dan $\Gamma \Proves !A$.
\end{prop}

\begin{proof}
  De !!{formula}~$!A$ op zichzelf is !!a{derivation} van~$!A$
  uit~$\Gamma$.
\end{proof}

\begin{prop}[Monotoniciteit]
\ollabel{prop:monotonicity}
Als $\Gamma \subseteq \Delta$ en $\Gamma \Proves !A$, dan $\Delta
\Proves !A$.
\end{prop}

\begin{proof}
Elke !!{derivation} van $!A$ uit $\Gamma$ is tevens !!a{derivation}
van $!A$ uit~$\Delta$.
\end{proof}

\begin{prop}[Transitiviteit]
\ollabel{prop:transitivity}
Als $\Gamma \Proves !A$ en $\{!A\} \cup \Delta \Proves
!B$, dan $\Gamma \cup \Delta \Proves !B$.
\end{prop}

\begin{proof}
  Veronderstel dat $\{!A\} \cup \Delta \Proves !B$. Dan bestaat er
  !!a{derivation} $!B_1$, \dots, $!B_l = !B$ uit~$\{!A\} \cup
  \Delta$. Sommige stappen in die !!{derivation} zijn correct op
  grond van een regel die naar een eerdere regel~$!B_i = !A$
  verwijst. Volgens de aanname bestaat er !!a{derivation} van~$!A$
  uit~$\Gamma$, dat wil zeggen !!a{derivation}~$!A_1$, \dots,
  $!A_k = !A$, waarin elke $!A_i$ een axioma, !!a{element}
  van~$\Gamma$ of correct volgens een afleidingsregel is. Beschouw nu
  de rij
  \[
  !A_1, \dots, !A_k = !A, !B_1, \dots, !B_l = !B.
  \]
  Dit is een correcte !!{derivation} van~$!B$ uit $\Gamma \cup
  \Delta$, want elke regel met $B_i = !A$ wordt nu gerechtvaardigd
  door dezelfde regel die~$!A_k = !A$ rechtvaardigt.
\end{proof}

Merk op dat hieruit in het bijzonder volgt dat, als $\Gamma \Proves
!A$ en $!A \Proves !B$, dan $\Gamma \Proves !B$. Hieruit volgt ook
dat, als $!A_1, \dots, !A_n \Proves !B$ en $\Gamma \Proves !A_i$
voor elke~$i$, dan $\Gamma \Proves !B$.

\begin{prop}
\ollabel{prop:incons}
$\Gamma$ is inconsistent dan en slechts dan als $\Gamma \Proves !A$
voor elke~$!A$.
\end{prop}

\begin{proof}
Oefening.
\end{proof}

\tagprob{FOL}
\begin{prob}
Bewijs \olref[fol][axd][ptn]{prop:incons}.
\end{prob}
\tagendprob

\tagprob{notFOL}
\begin{prob}
Bewijs \olref[pl][axd][ptn]{prop:incons}.
\end{prob}
\tagendprob

\begin{prop}[Compactheid]
\ollabel{prop:proves-compact}
  \begin{enumerate}
  \item Als $\Gamma \Proves !A$, dan bestaat er een eindige
    deelverzameling $\Gamma_0 \subseteq \Gamma$ waarvoor
    $\Gamma_0 \Proves !A$.
  \item Als elke eindige deelverzameling van~$\Gamma$ consistent is,
    dan is $\Gamma$ consistent.
  \end{enumerate}
\end{prop}

\begin{proof}
  \begin{enumerate}
    \item Als $\Gamma \Proves !A$, bestaat er een eindige rij
      !!{formula}s $!A_1$, \dots,~$!A_n$ waarvoor $!A \ident !A_n$,
      waarbij elke $!A_i$ ofwel een logisch axioma, ofwel !!a{element}
      van~$\Gamma$ is, ofwel door modus ponens uit eerdere
      !!{formula}s volgt. Laat $\Gamma_0$ bestaan uit die $!A_i$ die
      tot~$\Gamma$ behoren. De !!{derivation} is dan tevens
      !!a{derivation} uit~$\Gamma_0$, en dus $\Gamma_0 \Proves !A$.
    \item Dit is de contrapositie van~(1) voor het bijzondere geval
      $!A \ident \lfalse$.
\end{enumerate}
\end{proof}

\end{document}
